\section{Divisibility}\label{sec:divisibility}
%% FUNCTION: Shows how the EFX rate of round robin changes with the remainder q at three goods per
%% agent (m = 12 to 15 for n = 4). Divisibility decides when values are similar and fades when
%% values are very unequal, which is why the main figure shows three levels of dispersion. With
%% identical valuations and q > 0, round robin is never EFX (proven).

Divisibility asks whether the goods can be shared out evenly. With twelve goods and four people,
everyone picks three times; with thirteen, one person picks a fourth time. We measure this by the
remainder $q=m-nr$, the number of agents who receive one good more than the others. As defined in Section~\ref{sec:scenario}, $q=0$ means that the goods can be divided evenly. We keep $r=3$
goods per agent fixed and vary $q$ from $0$ to $3$, which means $m=12$ to $15$ goods for $n=4$
agents, at Spearman $0.48$ ($\rho=0.5$). Because the effect of divisibility depends on dispersion
(Section~\ref{sec:dispersion-rate}), Figure~\ref{fig:divisibility} shows three levels of
dispersion instead of the reference level alone.

\subsection{EFX Rate}\label{sec:divisibility-rate}

As Figure~\ref{fig:divisibility} shows, divisibility decides the outcome when the values of the
goods are similar, and hardly matters when they are not. At Gini $0.06$ ($\theta=0.1$), every one
of the $20{,}000$ instances with $m=12$ is EFX. With a single additional good, the EFX rate drops
to $0.499$, with $10{,}023$ failures; it is $0.466$ at $q=2$ and $0.605$ at $q=3$. At Gini $0.28$
($\theta=0.5$), the same pattern is weaker: $0.574$ at $q=0$ against $0.412$, $0.395$ and $0.475$.
At the reference level Gini $0.52$ ($\theta=1$), the advantage of divisibility has disappeared:
the EFX rate is $0.293$ at $q=0$ and $0.298$, $0.307$ and $0.342$ at $q=1$, $2$ and $3$. All of
these results are measured.

\begin{figure}[ht]
  \centering
  \includegraphics{fig44_divisibility_rate.pdf}
  \caption{EFX rate of Round Robin against the remainder $q$, for $n=4$ agents with $r=3$ goods
  per agent ($m=12$ to $15$) and Spearman $0.48$ ($\rho=0.5$), at three levels of dispersion.
  Each point rests on $R=20{,}000$ instances.}
  \label{fig:divisibility}
\end{figure}

Two further features are visible. Among the cells with $q>0$, the EFX rate is lowest in the
middle, at $q=2$, and highest at $q=3$. A plausible reason is the number of pairs in which an
agent with $r$ goods compares their bundle with a bundle of $r+1$ goods, which is $(n-q)\,q$ and
largest in the middle; this explains the shape only roughly and not its height, so we state it as
a conjecture. Moreover, at Gini $0.52$ the EFX rate rises slightly with $q$. Since the number of
goods rises with $q$, this may be the effect of more goods per agent from
Section~\ref{sec:goods}, which our design cannot separate from divisibility.

The pattern holds across the design. We compare the EFX rate at $q=0$ with the rate at $q=1$ for
the same $r$, for $r=2$ and $r=3$ and every number of agents, which gives $40$ comparisons for
each level of dispersion. Up to Gini $0.40$ ($\theta=0.75$), the divisible cell is ahead in all
$160$ comparisons, at every level of agreement. Below full agreement, its median lead is $0.44$
at Gini $0.06$ ($\theta=0.1$), $0.42$ at Gini $0.14$, $0.18$ at Gini $0.28$ and $0.06$ at Gini
$0.40$. At Gini $0.52$ ($\theta=1$), the divisible cell is ahead in $30$ of $40$ comparisons, with
a median lead of only $0.016$ below full agreement, and at Gini $0.71$ ($\theta=1.5$) in $13$ of
$40$, with a median of $-0.013$ (measured). The comparisons share their random seeds, so they
describe one pattern rather than independent confirmations.

\subsection{Identical Valuations without Divisibility}\label{sec:divisibility-identical}

Under full agreement, divisibility turns from an advantage into a strict condition.

\begin{proposition}\label{prop:identical}
Let all agents have identical valuations without ties, and let $q\geq1$. Then the Round Robin
allocation is not EFX, for every $n\geq2$.
\end{proposition}

\begin{proof}
Number the goods by their common value, so that rank $1$ is the most valuable. Since $q\geq1$,
agent $1$ receives the ranks $1,n+1,\dots,rn+1$, which are $r+1$ goods, and agent $n$ receives the
ranks $n,2n,\dots,rn$, which are $r$ goods. Agent $n$ removes from $A_1$ the least valuable good,
rank $rn+1$. The remaining goods have ranks $1,n+1,\dots,(r-1)n+1$, and each of them is more
valuable than the corresponding good $n,2n,\dots,rn$ of agent $n$, because $(k-1)n+1<kn$ for
$n\geq2$. Hence $v_n(A_1\setminus\{g\})>v_n(A_n)$ for this good $g$, and EFX fails.
\end{proof}

In our model, full agreement yields identical valuations without ties with probability one, so
Round Robin is then almost surely not EFX whenever $q\geq1$, at every level of dispersion. We
checked this in the stored cells: in all $180$ cells with $\rho=1$ and $q>0$, not a single
instance is EFX. The proof does not depend on our distribution of values. To our knowledge, the
statement has not been made explicitly before. It describes a weakness of Round Robin and not of
the problem: with identical valuations, an EFX allocation always exists
\parencite{plaut2020,branzei2025}. Chapter~\ref{ch:random} returns to this case and shows a simple
remedy.

\subsection{Severity of Failures}\label{sec:divisibility-severity}

Figure~\ref{fig:divisibility-survival} shows the survival curves of the four cells at the
reference level of dispersion. They lie close together, in line with the small differences in
the EFX rate, and their order follows $q$: a larger remainder gives a slightly higher curve. At
$\alpha=0.6$, the share of instances that reach at least this level is $0.651$, $0.680$, $0.704$
and $0.740$ for $q=0$ to $3$; the conditional $5\%$ quantile is $0.23$, $0.25$, $0.26$ and $0.28$,
and the median of $\astar$ among the instances that are not EFX is $0.60$, $0.63$, $0.65$ and
$0.67$. At the reference dispersion, a missing good to divide evenly therefore does not make
failures more severe. At low dispersion, the failures without divisibility are numerous but
narrow: at Gini $0.06$ ($\theta=0.1$), the conditional $5\%$ quantile lies between $0.92$ and
$0.93$ for $q=1$ to $3$, so even the worst failures miss EFX by less than ten percent (measured).

\begin{figure}[ht]
  \centering
  \includegraphics{fig44_divisibility_survival.pdf}
  \caption{Share of instances with $\astar\geq\alpha$ under Round Robin, one curve per remainder
  $q$, for $n=4$ agents with $r=3$ goods per agent ($m=12$ to $15$), Spearman $0.48$
  ($\rho=0.5$) and Gini $0.52$ ($\theta=1$). Each curve rests on $R=20{,}000$ instances. The point
  on each curve marks the $5\%$ quantile of $\astar$ among the instances that are not EFX.}
  \label{fig:divisibility-survival}
\end{figure}

\paragraph{Relation to literature.}
Divisibility plays a similar role for the existence of envy-free allocations.
\textcite{manurangsi2019} show that with independent values, envy-free allocations exist with
high probability when the number of goods is a multiple $rn$ with $r\geq2$, whereas without such
divisibility they are unlikely to exist even for considerably more goods. Our results show the
same dividing line for EFX under Round Robin, as long as values are not too dispersed, and
Proposition~\ref{prop:identical} makes it strict under full agreement. With identical valuations,
EFX existence is proven \parencite{plaut2020}, so every failure in these cells is a failure of
Round Robin. For $n=4$ and $m=12$ to $15$ below full agreement, existence is open and we assume it
(Section~\ref{sec:existence}).

\paragraph{Recommendation.}
When the goods are worth roughly the same, Round Robin should be used with a number of goods that
can be divided evenly among the people, and under full agreement it should never be used without
this divisibility; when values differ greatly, divisibility no longer matters. Having examined
what each person receives, the next section asks what changes when more people take part.
